\section{Conclusion: The Man in the Documents}

Strip away the two anachronisms---the plaster saint and the
grand inquisitor---and the documents leave a figure of unusual
consistency. The same man appears at every layer of the evidence:
the professor who taught himself Greek a lesson ahead of his
class; the preacher who gave up eloquence for bare points; the
controversialist who built a map of the whole disputed frontier
and was placed on the Index by his own side for drawing one
boundary too carefully; the courtier who fled Rome four days
after his consecration to reach his diocese, and whose own
account of Capua is a ledger of synods, visitations, and monthly
alms; the elector who prayed to be delivered from the papacy and
was; the old cardinal who asked the astronomers what the
telescope really showed, told Foscarini what would follow if a
true demonstration ever came, and signed a plain certificate so
that a friend under rumor could defend his name.

Historically, three judgments seem secure on the evidence
assembled here. First, Bellarmine was the most consequential
Catholic controversial theologian of his age: the
\work{Disputationes} organized the confessional debate for both
sides, and his primacy doctrine ran forward, by the official act's
own statement, to the Vatican Council. Second, his greatness was
administered almost entirely through institutions---chair,
congregations, legations, tribunal---so that his personal
holiness, which the process documented to the satisfaction of
every voting body that ever examined it, is inseparable in the
record from institutional acts, including coercive ones, whose
full particulars are often lost; the honest biography carries
both without netting them against each other. Third, the
three-century delay of his cause was the price of the second
fact: not doubt of the man, but the entanglement of his name with
a political theology that Catholic states would not see honored.
The Church that finally canonized him in 1930 and made him a
Doctor in 1931 was one that no longer possessed---or
wanted---the powers his century had argued about.

Of his sanctity the historian reports what the record shows: a
life whose austerity, prayer, and charity were attested by those
who lived beside it, examined by the most adversarial
canonization process on record, and found heroic by every
tribunal that voted on the question across three centuries. The
Church's judgment is the Church's act, stated in \work{Lux illa}
within its own object. The unresolved questions---his exact hand
in the Bruno sentence, the true course of the morning of
26~February 1616, the motives behind Capua, the full revision
history of the great work---are stated where they arise, and
remain open. He would perhaps not have minded. The man who wrote
\latin{De arte bene moriendi} in his annual retreat had already
made the only closing argument he cared about; the rest he left,
with his papers, to the tribunals of history---which are still
sitting.
